\subsection{Free Quotients and Splittings}

\begin{theorem} \label{thm:epimorphism-onto-free-module-splits} \hlt{Free Quotients Split}\\
Assume the Choice Principle \ref{princ:axiom-of-choice-analysis}. Let $A$ be a nonzero commutative ring with identity, and $M,F$ be $A$-modules.
\begin{enumerate}[label=(\roman*)]
\setlength{\itemsep}{0pt}
\item If $\sigma:M\twoheadrightarrow F$ and $F$ is free, then $(\exists\,\tau\in\operatorname{Hom}_A(F,M))[\sigma\tau=\operatorname{id}_F\land M=\ker\sigma\oplus\operatorname{im}\tau\land\operatorname{im}\tau\cong F]$.
\item If $S\leq M$ and $M/S$ is free, then $(\exists\,T\leq M)[M=S\oplus T\land T\cong M/S]$ and $M\cong S\oplus M/S$.
\item If $S\leq M$ and $M,S,M/S$ are free, then $\operatorname{rk}_A(M)=\operatorname{rk}_A(S)+\operatorname{rk}_A(M/S)$, where $+$ denotes cardinal addition.\\
If $\operatorname{rk}_A(M)\in\N_0$, then $\operatorname{rk}_A(M/S)=\operatorname{rk}_A(M)-\operatorname{rk}_A(S)$.
\end{enumerate}
\begin{proof}
\begin{enumerate}[label=(\roman*)]
\setlength{\itemsep}{0pt}
\item Fix $A$-basis $B$ of $F$. Relation $(\forall b\in B)\,\sigma^{-1}(\{b\})\neq\emptyset$ holds by surjectivity. Fix $(m_b)_{b\in B}\in\prod\limits_{b\in B}\sigma^{-1}(\{b\})$ by Choice Principle.\\
The assignment $b\mapsto m_b$ extends uniquely to $\tau\in\operatorname{Hom}_A(F,M)$ by Theorem \ref{thm:free-module-map-extension}.\\
The equality $\sigma\tau=\operatorname{id}_F$ holds by $(\forall b\in B)\,\sigma\tau(b)=\sigma(m_b)=b$ and uniqueness in Theorem \ref{thm:free-module-map-extension}.\\
Decomposition $M=\ker\sigma\oplus\operatorname{im}\tau$ holds by Theorem \ref{thm:module-one-sided-inverses-and-complements}(ii). Map $\tau$ is injective by $(\forall x\in F)[\tau(x)=0_M\Rightarrow x=\sigma\tau(x)=0_F]$.\\
Its corestriction $F\to\operatorname{im}\tau$ is an isomorphism by linearity and bijectivity.
\item Let $\pi:M\twoheadrightarrow M/S$ be the quotient map; $\ker\pi=S$ by Proposition \ref{prop:quotient-module-well-defined}(iii).\\
Fix a section $\tau\in\operatorname{Hom}_A(M/S,M)$ of $\pi$ by (i), and set $T\defeq\operatorname{im}\tau$. Relations $M=S\oplus T$, $T\cong M/S$ hold by (i) and $\ker\pi=S$.\\
Define $\Phi:S\oplus M/S\to M$ by $\Phi(s,q)\defeq s+\tau(q)$. Map $\Phi$ is $A$-linear by linearity of $\tau$, surjective by $M=S\oplus T$.\\
Identity $\ker\Phi=\{(0_M,0_{M/S})\}$ holds by $S\cap T=\{0_M\}$, injectivity of $\tau$. Map $\Phi$ injective by Proposition \ref{prop:hom-module-kernel-image}(iv), kernel identity.\\
It is an isomorphism by injectivity, surjectivity, and Definition \ref{def:module-homomorphism}(iv)(c). Relation $M\cong S\oplus M/S$ holds by $\Phi^{-1}:M\xrightarrow{\sim}S\oplus M/S$.
\item Fix $T\leq M$ with $M=S\oplus T$ and $\theta:T\xrightarrow{\sim}M/S$ by (ii). Fix $A$-bases $B_S$ of $S$ and $B_Q$ of $M/S$, and set $B_T\defeq\theta^{-1}(B_Q)$.\\
The set $B_T$ is an $A$-basis of $T$, and $\theta|_{B_T}:B_T\to B_Q$ is a bijection, by bijectivity and linearity of $\theta$.\\
Relation $B_S\cap B_T=\emptyset$ holds by $S\cap T=\{0_M\}$, $0_M\notin B_S\cup B_T$; $B_S\cup B_T$ spans $M$ by $M=S\oplus T$, spanning properties of $B_S,B_T$.\\
It is $A$-linearly independent by Proposition \ref{prop:module-direct-sum-characterizations}(ii) and the independence of $B_S,B_T$.\\
The union $B_S\cup B_T$ is an $A$-basis of $M$ by its spanning and independence.\\
The identities $\abs{B_S\cup B_T}=\abs{B_S}+\abs{B_T}$ and $\abs{B_T}=\abs{B_Q}$ hold by disjointness and the bijection $\theta|_{B_T}$.\\
Rank equality in (iii) holds by Definition \ref{def:rank-of-free-module}, identities. Under $\operatorname{rk}_A(M)\in\N_0$, both summand ranks lie in $\N_0$ by rank equality.\\
The formula $\operatorname{rk}_A(M/S)=\operatorname{rk}_A(M)-\operatorname{rk}_A(S)$ holds by subtraction in $\N_0$.
\end{enumerate}
\hlt{Structural Commentary}: Choice lifts a free target basis to a section; direct-sum bases then add ranks.\\
\hlt{Key Insight}: The section image is simultaneously a kernel complement and quotient model.
\end{proof}
\end{theorem}
